\begin{enumerate}\setlength{\itemsep}{1pt}
\item \textbf{2103.02708 Fig. 4, gpt-5.3-codex / codex}: High-mass normalisation agrees well (SSM \textasciitilde{}2.5e-5 pb and psi \textasciitilde{}1e-5 pb near 5 TeV, comparable to the reference theory curves), but at low mass the produced curves sit roughly a factor 2-4 high (psi \textasciitilde{}9e-2 pb at 1 TeV versus \textasciitilde{}2e-2 pb inferred from the reference), so the slope is slightly shallower than the paper's. The scan also starts near 600 GeV rather than covering the reference's full range down to \textasciitilde{}250 GeV.
\item \textbf{9909255 Fig. 2, gpt-5.3-codex / codex}: The produced scan covers 200-1200 GeV and therefore omits the Z pole spike near 91 GeV visible in the reference. Normalisation differs by roughly a factor of 3 at the 600 GeV resonance (\textasciitilde{}1.3e5 fb produced vs \textasciitilde{}4e4 fb in the reference) and is low by a factor of \textasciitilde{}2-3 in the 400 GeV dip (\textasciitilde{}9e2 fb vs \textasciitilde{}2e3 fb); the coarse 100 GeV step also under-resolves the resonance widths, so the 600 GeV peak height and the 1100 GeV structure are sampled rather than traced.
\item \textbf{1308.2209 Fig. 3, claude-opus-5}: Agent's 14 TeV σ(μ±N) agrees with the reference red curve to within \textasciitilde{}10–20\% across the mass range (e.g. \textasciitilde{}330 fb vs \textasciitilde{}380 fb at 300 GeV, \textasciitilde{}2.5 fb vs \textasciitilde{}3 fb at 1 TeV); the y-axis is quoted as σ [fb] at |V\_μN| = 1 rather than the reference's σ/|V\_ℓN|², which is numerically equivalent.
\item \textbf{2103.02708 Fig. 4, claude-opus-5}: Absolute normalisation tracks the reference well at high mass (SSM ≈2×10⁻⁵ pb near 5 TeV and Z'\_ψ ≈1.6×10⁻⁵ pb near 4.5 TeV, consistent with the reference crossing points to within \textasciitilde{}20-30\%). The produced figure uses a linear mass axis and extends down to 200 GeV, whereas the reference is logarithmic in mass and truncated at 10⁻¹ pb, so the low-mass part of the curves is off-scale in the reference. Two extra 'inclusive' (full Breit-Wigner, no mass window) curves are shown in addition to the resonant ones; they deviate from the resonant curves only above \textasciitilde{}3 TeV, by up to a factor \textasciitilde{}4 at 5.5 TeV for the SSM.
\item \textbf{1701.05379 Fig. 8, gemini-3.1-pro-preview / adk}: The produced y-axis is per-bin event fraction (\textasciitilde{}0.09 at the peak with \textasciitilde{}20 GeV bins) whereas the reference plots 1/sigma d(sigma)/dE\_T\textasciicircum{}miss in GeV\textasciicircum{}-1, so the two differ by a bin-width factor of \textasciitilde{}20; after that conversion the peak (\textasciitilde{}4.5e-3 GeV\textasciicircum{}-1) sits below the c1/c2/cW curves (\textasciitilde{}0.2) and the tail at 1 TeV (\textasciitilde{}1.3e-3 per bin, i.e. \textasciitilde{}6e-5 GeV\textasciicircum{}-1) falls between the c1/c2/cW (\textasciitilde{}3e-4) and c6/c8 (\textasciitilde{}3-6e-3) reference curves. The steeply-falling SM reference curve (which dies out by \textasciitilde{}250 GeV) is not reproduced, and no per-coefficient breakdown is shown.
\item \textbf{1308.2209 Fig. 3, gemini-3.1-pro-preview / gemini}: The produced 14 TeV curve is systematically \textasciitilde{}10-25\% below the reference in the mid-mass range (e.g. \textasciitilde{}1.2×10\textasciicircum{}3 fb vs \textasciitilde{}1.5×10\textasciicircum{}3 fb at m\_N = 200 GeV, \textasciitilde{}100 fb vs \textasciitilde{}115 fb at 400 GeV), while the endpoints agree well (\textasciitilde{}1.7×10\textasciicircum{}4 fb at 100 GeV and \textasciitilde{}2.5 fb vs \textasciitilde{}3 fb at 1000 GeV). The produced y-axis is σ in fb rather than the reference's σ/|V\_ℓN|², i.e. it corresponds to the mixing-independent normalisation |V\_ℓN|² = 1.
\item \textbf{2103.02708 Fig. 4, gemini-3.1-pro-preview / gemini}: Normalisation is not directly comparable since the reference axis is the ratio (σB)Z'/(σB)Z scaled by 1928 rather than a cross section in pb; taking the reference curves at face value, the produced Z'\_ψ curve sits high by roughly a factor of 3-5 at 1 TeV (\textasciitilde{}9×10⁻² pb vs \textasciitilde{}2×10⁻²) while being low by a similar factor near 5 TeV (\textasciitilde{}6×10⁻⁶ pb vs \textasciitilde{}2×10⁻⁵), i.e. the produced curves fall somewhat more steeply with mass.
\item \textbf{1308.2209 Fig. 3, gpt-5.3-codex / codex}: Produced 14 TeV values run slightly below the reference throughout: \textasciitilde{}1.5×10\textasciicircum{}4 vs \textasciitilde{}2×10\textasciicircum{}4 fb at 100 GeV, \textasciitilde{}1.2×10\textasciicircum{}3 vs \textasciitilde{}1.5×10\textasciicircum{}3 fb at 200 GeV, \textasciitilde{}100 vs \textasciitilde{}110 fb at 400 GeV, and \textasciitilde{}2.5 vs \textasciitilde{}3 fb at 1000 GeV — a roughly 10-25\% normalisation offset, consistent with the reference being normalised as σ/|V\_lN|\textasciicircum{}2 while the produced curve is a plain cross section at the chosen mixing.
\item \textbf{1605.02910 Fig. 1, gemini-3.1-pro-preview / gemini}: The reproduced LHC contour tips sit somewhat higher than the reference: g₁′ ≈ 0.17 vs ≈ 0.12 (2 TeV), ≈ 0.30 vs ≈ 0.27 (2.5 TeV), and ≈ 0.55 vs ≈ 0.65 (3 TeV, here lower); the lobes are also shifted \textasciitilde{}0.05-0.1 toward negative g̃ at 2 and 2.5 TeV, and the 3 TeV lobe is broader in g̃ (left edge near -0.65 vs -0.75 at the tip with a wider lower crossing spread).
\item \textbf{1701.05379 Fig. 8, gpt-5.3-codex / codex}: Only one coupling benchmark is shown, versus the seven curves (c1, c2, c6, c8, cB, cW, SM) of the reference; the produced curve tracks the c2/cW-like family, sitting within roughly a factor of two of them (e.g. \textasciitilde{}2.5e-2 vs \textasciitilde{}3e-2 at 200 GeV and \textasciitilde{}1.6e-3 vs \textasciitilde{}8e-4 at 600 GeV), and it does not reproduce the much harder c6/c8 spectra or the fast-falling SM and cB tails. The y-axis is labelled 'normalized events' rather than the reference's (1/σ)dσ/dE\_T\textasciicircum{}miss in GeV\textasciicircum{}-1, so the vertical scale agrees in magnitude but not strictly in units.
\item \textbf{9909255 Fig. 2, claude-opus-5}: Scan starts at 200 GeV, so the Z pole near 91 GeV and the low-√s rise visible in the reference are absent. The 100 GeV step size undersamples the resonance: the 600 GeV peak reaches \textasciitilde{}3.5x10\textasciicircum{}4 fb (reference peak \textasciitilde{}4x10\textasciicircum{}4 fb) but its true Breit-Wigner width is not resolved, and the second KK resonance near 1100 GeV appears only as a smooth rise to \textasciitilde{}2.7x10\textasciicircum{}4 fb rather than a distinct peak. The post-resonance minimum sits at \textasciitilde{}3x10\textasciicircum{}3 fb near 800-900 GeV versus \textasciitilde{}2.5x10\textasciicircum{}3 fb near 850 GeV for the reference blue curve.
\item \textbf{1605.02910 Fig. 1, gpt-5.3-codex / codex}: The LHC contour extent differs modestly from the reference: the maximum g1' reach is \textasciitilde{}0.17 vs \textasciitilde{}0.12 at M\_Z' = 2 TeV (\textasciitilde{}40\% high), \textasciitilde{}0.25 vs \textasciitilde{}0.27 at 2.5 TeV (close), and \textasciitilde{}0.48 vs \textasciitilde{}0.65 at 3 TeV (\textasciitilde{}25\% low); the 3 TeV contour's left edge also sits near g̃ ≈ −0.57 instead of ≈ −0.75, so the contour is somewhat narrower in g̃ at the highest mass. Lower-edge intercepts on g1' = 0 agree well (e.g. \textasciitilde{}−0.07 to +0.07 at 2 TeV vs \textasciitilde{}−0.09 to +0.05).
\item \textbf{9909255 Fig. 2, gemini-3.1-pro-preview / adk}: Sampled on a coarse 100 GeV grid (200-1200 GeV), so the resonances are rendered as single-point spikes rather than resolved Breit-Wigner shapes; the on-peak values (\textasciitilde{}1.6e6 fb at 600 GeV and >5e6 fb at 1100 GeV) overshoot the reference peak heights (\textasciitilde{}4e4 and \textasciitilde{}1e4 fb) by one to two orders of magnitude because the grid point sits exactly on the pole with no width smearing. The low-energy region below 200 GeV, including the Z pole and the SM continuum falloff, is not covered, and the off-peak continuum near 200-500 GeV (\textasciitilde{}1-3e3 fb) sits below the reference's \textasciitilde{}2e3-4e4 fb band.
\item \textbf{9909255 Fig. 2, gemini-3.1-pro-preview / gemini}: The scan covers only √s = 400–1200 GeV (points every 100 GeV), so the reference's low-energy region (Z pole near 90 GeV and the falling SM QED tail down to the \textasciitilde{}300 GeV minimum) and the region up to 1500 GeV are absent; the reference's family of curves for different parameter choices is represented by a single curve. The lowest scan point (≈400 GeV) falls below \textasciitilde{}4×10\textasciicircum{}2 fb and runs off the bottom of the frame, whereas the reference curves are \textasciitilde{}1–2×10\textasciicircum{}3 fb there, i.e. the SM-dominated low-energy behaviour is not reproduced. At the resonance the produced peak is ≈3.3×10\textasciicircum{}4 fb versus ≈4×10\textasciicircum{}4 fb in the reference, and the 800 GeV dip (≈2.7×10\textasciicircum{}3 fb) and 1200 GeV value (≈2.6×10\textasciicircum{}4 fb) agree with the reference blue curve to within roughly 10–20\%; the coarse 100 GeV grid also means the resonance peak height is only sampled, not resolved.
\item \textbf{2104.05720 Fig. 11, claude-opus-5}: In the left panel the produced m\_LQ = 1 TeV curve peaks slightly further forward (eta \textasciitilde{} 1.3-1.5 vs \textasciitilde{} 1.1-1.4 in the reference) and stays \textasciitilde{}30-60\% higher in the tail (\textasciitilde{}0.04 vs \textasciitilde{}0.02-0.03 at eta \textasciitilde{} 2.3-2.5); the reference plots eta\_mu whereas the produced figure labels the axis simply eta.
\item \textbf{2104.05720 Fig. 11, gpt-5.3-codex / codex}: SM baseline reproduces the reference well (first bin \textasciitilde{}0.075, peak \textasciitilde{}0.150-0.155 at eta\textasciitilde{}0.2, falling to \textasciitilde{}0.01 at eta\textasciitilde{}2.4). Deviation is in the BSM curves: reference mLQ=1 TeV, betaL32=1 peaks at \textasciitilde{}0.13 near eta\textasciitilde{}1.4 and starts at \textasciitilde{}0.02 in the first bin, whereas the produced curve peaks at \textasciitilde{}0.151 at eta\textasciitilde{}0.2 and starts at \textasciitilde{}0.075, i.e. matches the SM within \textasciitilde{}1-2\% everywhere; the mLQ=10 TeV curve agrees with the reference (slightly above SM in the first two bins, \textasciitilde{}0.155 vs \textasciitilde{}0.150).
\item \textbf{1308.2209 Fig. 3, gpt-5.5 / adk}: At √s = 14 TeV the produced cross sections track the reference pp→Nℓ± curve closely over the full mass range (e.g. \textasciitilde{}1.5×10\textasciicircum{}4 fb vs \textasciitilde{}2×10\textasciicircum{}4 fb at 100 GeV, \textasciitilde{}1.2×10\textasciicircum{}3 vs \textasciitilde{}1.5×10\textasciicircum{}3 fb at 200 GeV, \textasciitilde{}105 vs \textasciitilde{}110 fb at 400 GeV, \textasciitilde{}2.6 vs \textasciitilde{}3 fb at 1000 GeV), i.e. agreement within roughly 10–25\%, with the produced values sitting slightly low at the lightest masses. The reference normalises to σ/|V\_ℓN|\textasciicircum{}2 while the produced figure fixes |V\_μN| = 1, which is equivalent.
\item \textbf{1701.05379 Fig. 8, gpt-5.5 / adk}: The produced histogram is normalized to unit total area per bin (dimensionless 'normalized events') rather than the reference's 1/sigma dsigma/dMET in GeV\textasciicircum{}-1, so absolute y-values are not directly comparable. The shape most closely tracks the reference c\_1/c\_2/c\_W curves; the harder c\_6/c\_8 curves (which peak around 150-250 GeV and fall far more slowly) are not represented, and the produced tail falls somewhat more slowly than the SM/c\_B curves.
\item \textbf{2104.05720 Fig. 11, gemini-3.1-pro-preview / gemini}: In the beta\_L\textasciicircum{}32 = 1.0 panel the reproduced m\_LQ = 1 TeV curve peaks slightly earlier (|eta| \textasciitilde{} 1.1-1.3 at \textasciitilde{}0.118) than the reference (\textasciitilde{}1.3-1.5 at \textasciitilde{}0.13) and stays higher in the tail (\textasciitilde{}0.04 vs \textasciitilde{}0.02 near |eta| \textasciitilde{} 2.4). In the beta = 0.1 panel the reproduced 1 TeV curve is slightly more peaked at low |eta| (\textasciitilde{}0.17 vs \textasciitilde{}0.175 leading bin, with a marginally deeper dip near |eta| \textasciitilde{} 1.2). The reference plots signed eta\_mu while the reproduction uses |eta|, which is equivalent given the symmetric range shown.
\item \textbf{1811.07920 Fig. 3, gpt-5.3-codex / codex}: The produced plot omits the reference's auxiliary grey EFT dash-dotted line and the 150 fb⁻¹ / 3 ab⁻¹ projection curves, and carries no in-panel RH/LH/EFT text labels. The black 2σ contour nearly coincides with the upper edge of the red RH band over the whole mass range, whereas in the reference the grey-region boundary sits slightly above the RH band at low mass; endpoint values (≈0.88 at 0.8 TeV, ≈3.75 at 5 TeV for the contour; ≈0.3→1.95 for LH) are within a few percent of the reference.
\item \textbf{2104.05720 Fig. 12, claude-opus-5}: Limits are somewhat weaker than the reference at large mass: the 3 TeV 5-sigma curve reaches beta\_L\textasciicircum{}32 \textasciitilde{} 2 at m\_LQ = 55 TeV versus \textasciitilde{}1 in the reference, and the 14 TeV 5-sigma curve \textasciitilde{}0.45 versus \textasciitilde{}0.3. At low mass the agreement is good for the discovery curves (\textasciitilde{}0.05 and \textasciitilde{}0.022 at 1 TeV vs \textasciitilde{}0.04 and \textasciitilde{}0.02), but the 14 TeV 95\% CL exclusion is notably less constraining at 1 TeV (\textasciitilde{}0.019 vs \textasciitilde{}0.006 in the reference), so the reproduced exclusion/discovery separation at 14 TeV is narrower than in the paper.
\item \textbf{1701.05379 Fig. 8, gemini-3.1-pro-preview / adk}: Only one curve is drawn, so no SM-vs-operator comparison is possible; its tail (\textasciitilde{}1e-3 at 1000 GeV, falling \textasciitilde{}2 decades from the peak) is much harder than the reference SM curve (which dies below 1e-4 by \textasciitilde{}250 GeV) and closest to the c6 curve, while the reference c8 curve is roughly a factor of \textasciitilde{}5 higher at 1 TeV. The y-axis is plotted as per-bin 'normalized events' rather than the reference's (1/sigma) dsigma/dE\_T [GeV\textasciicircum{}-1], so absolute vertical values differ by the bin width (peak \textasciitilde{}0.12 vs \textasciitilde{}0.2-0.5 in the reference).
\item \textbf{2103.02708 Fig. 4, gemini-3.1-pro-preview / adk}: The reference y-axis is the scaled ratio (σB)Z'/(σB)Z × 1928 rather than an absolute cross section in pb, so the normalisations are not directly comparable. More notably, the relative ordering of the two benchmarks is inverted: in the reference the Z'\_SSM (dotted) curve lies above Z'\_psi over the whole mass range (σ\_SSM ≈ 3× σ\_psi), while the produced figure places Z'\_extSSM roughly a factor 2–3 below Z'\_psi. The produced Z'\_psi curve also reaches \textasciitilde{}5×10\textasciicircum{}-6 pb only near 5.1 TeV, i.e. the two curves cross the low-cross-section region a few hundred GeV apart from the reference.
\item \textbf{2104.05720 Fig. 12, gemini-3.1-pro-preview / gemini}: Produced contours sit \textasciitilde{}20-30\% above the reference at high mass (3 TeV discovery reaches \textasciitilde{}1.8 vs \textasciitilde{}1.1 at 50 TeV; 14 TeV discovery \textasciitilde{}0.5 vs \textasciitilde{}0.3). The 14 TeV 95\% CL curve is too close to its 5-sigma counterpart at low mass (\textasciitilde{}0.016 at 1 TeV vs \textasciitilde{}0.006 in the reference), so the exclusion/discovery gap is compressed. Minor non-physical kinks appear near m\_LQ = 1.5-2 TeV in the 14 TeV curves, likely from a coarse mass grid.
\item \textbf{9909255 Fig. 2, gemini-3.1-pro-preview / adk}: The scan uses \textasciitilde{}100 GeV steps, so the resonances appear as triangular spikes rather than resolved Breit-Wigner peaks: the 600 GeV point reaches >4×10\textasciicircum{}6 fb and the 1100 GeV point ≈2×10\textasciicircum{}6 fb, roughly two orders of magnitude above the reference peak heights (\textasciitilde{}4×10\textasciicircum{}4 and \textasciitilde{}10\textasciicircum{}4 fb), while the inter-resonance minimum near 800 GeV falls below the plotted range instead of the reference's \textasciitilde{}1.5×10\textasciicircum{}2 fb floor. The plotted range starts near 420 GeV, so the low-√s region including the Z pole and the SM continuum fall-off (√s ≲ 400 GeV in the reference) is not covered, and the curve stops at 1200 GeV rather than 1500 GeV.
\item \textbf{1605.02910 Fig. 1, gpt-5.5 / adk}: Panel (c) (M\_Z' = 3 TeV) reproduces a somewhat smaller contour than the reference: the tip reaches g1' ≈ 0.48 at g̃ ≈ −0.55 versus ≈ 0.65 at g̃ ≈ −0.72 in the paper (\textasciitilde{}25-30\% low in g1' and in leftward extent), although the right-hand g1' = 0 intercept (g̃ ≈ 0.18) agrees. Panels (a) and (b) agree to within a few percent (tip ≈ (−0.15, 0.12) vs (−0.13, 0.12); ≈ (−0.28, 0.245) vs (−0.30, 0.265)).
\item \textbf{2103.02708 Fig. 4, gpt-5.5 / adk}: Normalisation runs somewhat high relative to the reference theory curves: the produced Z'\_SSM cross section is \textasciitilde{}1×10⁻⁴ pb at 4.3 TeV versus \textasciitilde{}3×10⁻⁵ pb in the reference (factor \textasciitilde{}3), while the Z'\_ψ curve agrees to within roughly a factor of 1.5 (\textasciitilde{}4×10⁻⁵ pb vs \textasciitilde{}3×10⁻⁵ pb near 4 TeV). The produced curve stops at the point where it leaves the plotted y-range (\textasciitilde{}5.2 TeV for Z'\_ψ) rather than extending to 5.5 TeV.
\item \textbf{1308.2209 Fig. 3, gemini-3.1-pro-preview / adk}: The produced 14 TeV curve sits \textasciitilde{}20–30\% below the reference throughout: \textasciitilde{}1.6×10\textasciicircum{}4 fb vs \textasciitilde{}2×10\textasciicircum{}4 fb at m\_N = 100 GeV, \textasciitilde{}300 fb vs \textasciitilde{}400 fb at 300 GeV, and \textasciitilde{}2.5 fb vs \textasciitilde{}3 fb at 1000 GeV. The reference normalises to σ/|V\_lN|\textasciicircum{}2 while the produced figure plots σ in fb at an implicit mixing value, so the offset may be partly a normalisation convention rather than a physics difference.
\item \textbf{2104.05720 Fig. 11, gemini-3.1-pro-preview / adk}: The SM-background (and 10 TeV) histogram exceeds the 0.20 y-axis limit and is clipped over 0.1 < |eta| < 1.1, while the reference SM distribution peaks at \textasciitilde{}0.15 near eta \textasciitilde{} 0.2 and falls monotonically; the 1 TeV signal peaks at |eta| \textasciitilde{} 1.0 instead of \textasciitilde{}1.3-1.5, and its high-|eta| tail (\textasciitilde{}0.06 at 2.5) is about 4x the reference value.
\item \textbf{1605.02910 Fig. 1, gemini-3.1-pro-preview / adk}: The 3 TeV LHC contour is noticeably smaller than the reference: its apex reaches g₁' ≈ 0.47 at g̃ ≈ −0.55 versus ≈ 0.65 at g̃ ≈ −0.65, and its left-hand zero crossing sits at g̃ ≈ −0.2 instead of ≈ −0.25. The 2.5 TeV contour is also slightly compressed (apex g₁' ≈ 0.24 vs ≈ 0.27, and it is thinner in the g̃ direction near the apex), while the 2 TeV panel agrees to within \textasciitilde{}10\% (apex g₁' ≈ 0.13 vs ≈ 0.12). The x-axis label is mislabelled 'ildeg' rather than g̃.
\item \textbf{9909255 Fig. 2, gpt-5.5 / adk}: The scan covers 200-1200 GeV in 100 GeV steps, so the Z pole at 91 GeV and the region above 1200 GeV present in the reference are not sampled, and the coarse binning under-resolves the narrow 600 GeV resonance (peak reaches \textasciitilde{}3x10\textasciicircum{}4 fb vs \textasciitilde{}4x10\textasciicircum{}4 fb for the reference blue curve). Only a single coupling choice is shown, whereas the reference overlays six curves; the near-1100 GeV second KK resonance is therefore only visible as a rising edge rather than a resolved peak.
\item \textbf{2104.05720 Fig. 11, gpt-5.5 / adk}: In the beta\_L\textasciicircum{}32 = 1.0 panel the m\_LQ = 1 TeV curve falls off more slowly at large pseudorapidity than in the reference: the produced distribution is \textasciitilde{}0.04 in the last bin (|eta| \textasciitilde{} 2.3-2.5) versus \textasciitilde{}0.01-0.02 in the reference, with a correspondingly broader peak (\textasciitilde{}1.1-1.5 at 0.117 versus a sharper peak near 1.4-1.5 at \textasciitilde{}0.13). The axis is labelled |eta\_b| (b-quark) whereas the reference labels it eta\_mu, though the plotted range and content are the same.
\item \textbf{2103.02708 Fig. 4, gemini-3.1-pro-preview / adk}: Normalisation runs roughly 2-4x above the reference theory curves (e.g. Z'\_psi at 2 TeV: \textasciitilde{}4e-3 pb produced vs \textasciitilde{}1e-3 in the reference, where the y-axis is (sigma*B)\_\{Z'\}/(sigma*B)\_Z x 1928); the high-mass tails converge better (Z'\_SSM at \textasciitilde{}4.7 TeV: \textasciitilde{}7e-5 vs \textasciitilde{}1e-4). The dotted curve is labelled Z'\_extSSM rather than the reference's Z'\_SSM.
\item \textbf{2104.05720 Fig. 12, gpt-5.5 / adk}: The 3 TeV curves agree closely with the reference (β\_L\textasciicircum{}32 ≈ 0.04–0.05 at m\_LQ = 1 TeV, ≈ 1–2 at 50 TeV). The 14 TeV exclusion curve is roughly 2–3× weaker than the reference at low mass (≈0.016 vs ≈0.006 at m\_LQ = 1 TeV) and converges to within \textasciitilde{}20–30\% above \textasciitilde{}10 TeV; the produced 14 TeV curves also show a small non-monotonic kink near m\_LQ ≈ 1.5–2 TeV (likely statistical noise in the scan) that is absent from the smooth reference contours.
\item \textbf{2104.05720 Fig. 12, gemini-3.1-pro-preview / adk}: Sensitivity is \textasciitilde{}20-60\% weaker than the reference at low mass (3 TeV 5σ: 0.048 vs \textasciitilde{}0.04 at m\_LQ = 1 TeV) and roughly a factor 1.5-2 weaker at the high-mass end (3 TeV 5σ reaches \textasciitilde{}2 vs \textasciitilde{}1.2 near 50 TeV). The 14 TeV 95\% CL curve is the largest outlier: \textasciitilde{}0.017 at 1 TeV versus \textasciitilde{}0.006 in the reference, so the separation between the 14 TeV 5σ and 95\% CL contours is much narrower than in the paper. The curves also show visible kinks (coarse mass grid) where the reference is smooth, and the reference's non-monotonic dip near 1.5-2 TeV in the 14 TeV curves is an artefact not present in the paper.
\item \textbf{2005.06475 Fig. 2, gpt-5.5 / adk}: The reproduced signal peak reaches \textasciitilde{}1.8x10\textasciicircum{}2 events/bin/100 fb\textasciicircum{}-1 versus \textasciitilde{}4 in the reference black LQ curve, while the low-mass end (\textasciitilde{}0.04 at 1 TeV) is close to the reference (\textasciitilde{}0.02), i.e. the peak-to-shoulder contrast is roughly an order of magnitude larger, consistent with the reference curve being drawn after the full event selection (and possibly a different branching-ratio/normalisation convention). The peak position (\textasciitilde{}3 TeV) and the shape of the falling high-mass tail agree.
\item \textbf{1811.07920 Fig. 3, gpt-5.5 / adk}: The reproduced 2σ exclusion boundary sits somewhat above the reference one at high mass (≈2.45 at 5 TeV vs ≈3.5 for the reference's boundary between the grey and white regions), so the excluded region is smaller; the reference exclusion edge tracks near the RH band up to \textasciitilde{}1.5 TeV and then runs parallel above it, whereas the reproduced contour bends below the RH band above \textasciitilde{}1.3 TeV, crossing it. Band widths are slightly wider than the reference, and the produced contour shows mild kinking from a coarse mass grid.
\item \textbf{2005.06475 Fig. 2, claude-opus-4-6}: TBD (handed over). Stopped by the operator after 6.4 h without a figure: model, validation and the 100k-event LUXlep parton sample were done; two Pythia8+Delphes launches failed at upload with 'No space left on device' (27 GB Delphes ROOT against the job's 10 GB packing space); a third launch with post-run cleanup was running when stopped.
\item \textbf{1605.02910 Fig. 1, claude-opus-4-6}: The 2 TeV and 2.5 TeV contours match the reference to within about 10\% (tips at (-0.155, 0.13) vs (-0.14, 0.12) and (-0.28, 0.245) vs (-0.30, 0.265)). At 3 TeV the base on g1'=0 matches but the tip is at (-0.55, 0.48) vs (-0.74, 0.65), about 25\% smaller in both g1' and |gt|. An Opus 5 run and the Opus 4.6 attempt 3 gave the same 3 TeV tip, whereas the Codex/gpt-5.5 run of 2026-09-24 reproduced the paper's tip (-0.72, 0.65): the shortfall is an analysis choice made in the Claude runs, not a limitation of the prompt (earlier wording 'systematic to the prompt' is withdrawn).
\item \textbf{1811.07920 Fig. 3, claude-opus-4-6}: Attempt 1 generated all 18 Delphes runs (4.0 h, 106 USD; four 9-run launches were lost to the 10 GB job limit before the agent split them), then the orchestrator ended its turn 'waiting for event file downloads' and the headless session terminated without the analysis stage. A continuation session in the same workspace produced the figure in 1.7 h / 29.5 USD with the reference's shape but a 20-35\% weaker exclusion. Counted as a failure (infrastructure).
\item \textbf{9909255 Fig. 2, claude-opus-4-6}: The 11-point sqrt(s) grid (200-1200 GeV, 100 GeV steps) and the single coupling point are what the benchmark prompt specifies, so the reference's Z-pole region, its 1.2-1.5 TeV range and its other coupling curves are outside the task. The grid under-resolves the 600 GeV KK resonance: produced peak about 1.5e4 fb against about 3e4 fb in the reference; off-peak values follow the reference curve.
\item \textbf{2104.05720 Fig. 11, claude-opus-4-6}: The orchestrator was not triggered: the main session ran the pipeline itself (0 sub-agent calls); the result is correct.
\item \textbf{2104.05720 Fig. 12, claude-opus-4-6}: The produced limits are weaker than the reference throughout: the 14 TeV 95\% CL curve at m\_LQ = 1 TeV is about 0.015 against about 0.006 (roughly 2.5x weaker); at 70 TeV the 14 TeV curves are about 0.25 (exclusion) and 0.4 (discovery) against about 0.18 and 0.25. The 3 TeV curves agree within 20-30\% at low mass but are about 1.6x weaker at 70 TeV for the exclusion.
\item \textbf{1701.05379 Fig. 8, claude-sonnet-5}: Cross section 0.04909 pb instead of the 0.04135 pb obtained by Opus 5 and Opus 4.8, from looser generation-level eta cuts (5 instead of 2.5); the normalization removes the difference.
\item \textbf{1605.02910 Fig. 1, claude-opus-5}: At M\_Z' = 3 TeV the produced contour tip is at g1' about 0.48 (gt about -0.55) against about 0.65 (gt about -0.73) in the reference, about 25\% lower; the g1' = 0 intercepts (about +-0.17) agree. 2 TeV tip about 0.135 against 0.12; 2.5 TeV tip about 0.245 against 0.265. Same 3 TeV contour as the Opus 4.6 runs; see the 1605.02910 Opus 4.6 attempt-1 note.
\item \textbf{1308.2209 Fig. 3, claude-opus-4-6}: Normalisation convention matches (|V\_muN| = 1 vs the reference's sigma/|V\_lN|\textasciicircum{}2). Point-by-point the 14 TeV curve sits \textasciitilde{}20-30\% below the reference at the low-mass end (1.4e4 vs \textasciitilde{}2e4 fb at m\_N = 100 GeV, 1.2e3 vs \textasciitilde{}1.4e3 fb at 200 GeV) and agrees to within \textasciitilde{}10-20\% above 400 GeV; the produced curve uses CTEQ6L, and no PDF/scale uncertainty band is shown.
\item \textbf{1605.02910 Fig. 1, claude-opus-4-6}: Contour tips are systematically slightly lower/narrower than the reference: panel (a) peaks at g₁′≈0.13 at g̃≈−0.15 vs ≈0.12 at ≈−0.13 (close); panel (b) peaks at ≈0.24 vs ≈0.27 and its left tail closes near g̃≈−0.3 vs ≈−0.32; panel (c) is the largest deviation, peaking at g₁′≈0.47 at g̃≈−0.55 versus ≈0.65 at ≈−0.7 in the reference, i.e. the 3 TeV exclusion region is roughly 25-30\% smaller in the vertical extent. The g̃-axis zero crossings (≈0.05, ≈0.1, ≈0.2 for the three masses) agree well with the reference.
\item \textbf{1605.02910 Fig. 1, claude-opus-4-6}: The contour tips are close for the two lighter masses (g1' ≈ 0.13 vs 0.117 at 2 TeV; 0.24 vs 0.265 at 2.5 TeV) but the 3 TeV panel is noticeably smaller: peak g1' ≈ 0.47 vs ≈ 0.65 in the reference, and the left edge reaches g̃ ≈ −0.57 instead of ≈ −0.78, i.e. a \textasciitilde{}25-30\% under-coverage of the excluded region at the highest mass. Zero-crossings on the g̃ axis agree to within \textasciitilde{}0.05 in all panels. The session ended with 'OAuth session expired and could not be refreshed' after the figure and the step-4 sidecar were written (only the execution summary is missing); the figure is judged as produced.
\item \textbf{1605.02910 Fig. 1, gpt-5.5 / codex}: Panel (a) and (c) contours match the reference almost exactly (tip at g₁'≈0.12 at g̃≈−0.13, and g₁'≈0.65 at g̃≈−0.72 respectively); panel (b) reproduces the tip height (g₁'≈0.26) and right-hand crossing (g̃≈+0.09) but the lobe is somewhat narrower/more pinched than the reference, whose contour is fuller in the 0.1<g₁'<0.25 region.
\item \textbf{1701.05379 Fig. 8, claude-opus-4-6}: The reference shows seven overlaid curves (c1, c2, c6, c8, cB, cW and SM) while the reproduction shows a single distribution; its shape tracks the c1/c2/cW-type curves (peak 40-80 GeV, \textasciitilde{}3e-4 at 1 TeV) rather than the much softer SM curve, which dies out by \textasciitilde{}250 GeV. The y-axis is a per-bin fraction ('normalized events') rather than the reference's density (1/sigma)(dsigma/dE\_T\textasciicircum{}miss) in GeV\textasciicircum{}-1, so the absolute level differs by roughly the 20 GeV bin width (peak \textasciitilde{}0.18 per bin vs \textasciitilde{}0.2 GeV\textasciicircum{}-1); the intermediate region near 200 GeV sits about a factor of 2 above the c1 curve and below c2.
\item \textbf{1701.05379 Fig. 8, claude-opus-4-6}: The produced histogram plots normalized events per bin rather than the reference's (1/σ)dσ/dE\_T density in GeV\textasciicircum{}-1, so absolute ordinate values differ by the bin width (peak 0.18 per 25 GeV bin vs \textasciitilde{}0.2-0.25 GeV\textasciicircum{}-1 in the reference). The tail at 1 TeV lands at \textasciitilde{}3x10\textasciicircum{}-4, consistent with the reference's c\_1/c\_2/c\_W curves but \textasciitilde{}20x below its c\_6/c\_8 curves; only one benchmark coefficient is shown, and no SM curve is overlaid.
\item \textbf{1701.05379 Fig. 8, gpt-5.5 / codex}: The produced histogram falls more slowly than the reference's c\_1/c\_2/c\_W/c\_B curves and faster than c\_6/c\_8: at 1000 GeV it sits at \textasciitilde{}2e-4, versus \textasciitilde{}3e-4 (c\_1/c\_2/c\_W), \textasciitilde{}3e-3 (c\_6) and \textasciitilde{}6e-3 (c\_8) in the reference. The peak height is \textasciitilde{}0.18 vs \textasciitilde{}0.2-0.25 for most reference curves, and the low-E\_T rise is smoother (finer binning) than the reference's sharp SM-like spike at \textasciitilde{}50 GeV. The produced curve is a per-bin normalized fraction rather than the reference's 1/sigma dsigma/dE\_T density, so absolute y-values differ by the bin width.
\item \textbf{1811.07920 Fig. 3, gpt-5.5 / codex}: The LH and RH R\_D(*) bands agree closely with the reference (LH from \textasciitilde{}0.3 at 0.8 TeV to \textasciitilde{}1.95 at 5 TeV; RH from \textasciitilde{}0.65 to \textasciitilde{}4). The exclusion boundary is somewhat weaker/steeper than the reference: it starts near 1.07 at 0.8 TeV (vs \textasciitilde{}0.8 in the reference, where the boundary essentially tracks the RH band at low mass) and reaches 4.0 near M\_U1 ≈ 3.9 TeV, whereas the reference excluded region hugs the RH band up to \textasciitilde{}5 TeV. The reference's EFT-validity line and the 150 fb\textasciicircum{}-1 and 3 ab\textasciicircum{}-1 projection curves are not drawn.
\item \textbf{2005.06475 Fig. 2, gpt-5.5 / codex}: The signal peak sits at \textasciitilde{}1.5 events/bin versus \textasciitilde{}3 in the reference (roughly a factor 2 lower normalisation), and the produced distribution terminates at \textasciitilde{}3.8 TeV whereas the reference LQ curve extends to 5 TeV with a low tail; the low-mass shoulder is also \textasciitilde{}2-5x below the reference near 1-2 TeV, and the x-axis is drawn linearly rather than logarithmically.
\item \textbf{2103.02708 Fig. 4, claude-opus-4-6}: The produced curves sit high relative to the reference theory lines: σB(Z'\_SSM) ≈ 2.6×10⁻¹ pb at 1 TeV versus ≈2–3×10⁻² pb in the reference (factor \textasciitilde{}10), narrowing to ≈4×10⁻⁵ pb versus ≈2×10⁻⁵ pb at 5 TeV (factor \textasciitilde{}2), i.e. the falling slope is somewhat shallower than the reference. The SSM/ψ ratio is roughly constant at \textasciitilde{}3 across the mass range, close to the reference ordering and spacing.
\item \textbf{2103.02708 Fig. 4, claude-opus-4-6}: The reference y-axis is the normalized ratio ((σB)Z'/(σB)Z)×1928 pb rather than a raw cross section; agreement is good at high mass (both \textasciitilde{}1e-5 pb for Z'\_SSM near 5 TeV) but the produced curves sit roughly a factor of a few higher below \textasciitilde{}1.5 TeV (produced Z'\_ψ ≈ 0.1 pb at 1 TeV vs \textasciitilde{}2e-2 in the reference), and the Z'\_ψ curve stops at \textasciitilde{}5.0 TeV instead of extending to 5.5 TeV.
\item \textbf{2104.05720 Fig. 11, claude-opus-4-6}: Produced curves deviate from SM by at most \textasciitilde{}2-3\% in any bin for both m\_LQ = 1 TeV and 10 TeV at β\_L\textasciicircum{}32 = 1, versus the reference where the 1 TeV curve is shifted by O(100\%) at |η| > 1.2 and suppressed by a factor \textasciitilde{}3 in the first bin. The 10 TeV curve and the β\_L\textasciicircum{}32 = 0.1 panel agree with the reference at the few-percent level. Normalisation convention also differs (produced uses 1/N dN/d|η| peaking at \textasciitilde{}1.5; reference plots per-bin normalized fractions peaking at \textasciitilde{}0.15), which is a bin-width factor, not a shape difference.
\item \textbf{2104.05720 Fig. 11, claude-opus-4-6}: In the left panel (β\_L\textasciicircum{}32 = 1) the 1 TeV curve peaks at η ≈ 1.3–1.5 with a normalized height of ≈0.118, slightly broader and \textasciitilde{}20\% lower than the reference peak of ≈0.145 at η ≈ 1.1–1.3, and it stays somewhat above the reference at large η (≈0.04 vs ≈0.01 in the last bin). The right panel agrees to within a few percent bin-by-bin. Axis is labelled η rather than η\_μ.
\item \textbf{2104.05720 Fig. 12, claude-opus-4-6}: The produced x-range stops near \textasciitilde{}50 TeV like the reference but the low-mass end is largely consistent (3 TeV 95\% CL \textasciitilde{}0.035 vs \textasciitilde{}0.04, 14 TeV 5-sigma \textasciitilde{}0.022 vs \textasciitilde{}0.02 at m\_LQ = 1 TeV); the 14 TeV 95\% CL curve sits \textasciitilde{}2-3x higher than the reference at low mass (\textasciitilde{}0.017 vs \textasciitilde{}0.006 at 1 TeV) and shows a spurious bump/dip structure around m\_LQ = 1-2 TeV that is absent in the smooth reference contours; at the high-mass end the produced 3 TeV curves rise somewhat faster (5-sigma reaching \textasciitilde{}2 vs \textasciitilde{}1 near 50 TeV), and the gap between 95\% CL and 5-sigma is narrower for the 14 TeV pair than in the reference.
\item \textbf{2104.05720 Fig. 12, claude-opus-4-6}: Overall normalisation is close (3 TeV 5σ ≈ 0.05 at 1 TeV and ≈1–2 at \textasciitilde{}50 TeV, matching the reference within tens of percent), but the 14 TeV 95\% CL curve sits high at low mass (≈0.016 at m\_LQ = 1 TeV versus ≈0.006 in the reference), so the exclusion–discovery gap at 14 TeV is compressed relative to the paper; the produced 14 TeV curves also show small non-physical kinks near m\_LQ ≈ 1.5–3 TeV where the reference is smooth.
\item \textbf{9909255 Fig. 2, claude-opus-4-6}: Sampled on a coarse 100 GeV grid (200-1200 GeV), so the 600 GeV resonance is sampled rather than resolved: the peak reaches \textasciitilde{}3.5e4 fb versus \textasciitilde{}4e4 fb in the reference, and the interpolated line makes the resonance triangular rather than Breit-Wigner shaped. The scan omits the Z-pole region below 200 GeV shown in the reference and stops at 1200 GeV rather than 1500 GeV, so any second KK resonance near 1100 GeV is not probed. The low-energy end is somewhat high (\textasciitilde{}3e3 fb at 200 GeV versus \textasciitilde{}2e3 fb for the corresponding reference curve), and the 800 GeV dip minimum is shallower/shifted by up to one grid spacing.
\item \textbf{9909255 Fig. 2, claude-opus-4-6}: Sampled on a coarse 100 GeV grid from 200–1200 GeV, so the Z pole near 91 GeV and the 1500 GeV endpoint of the reference are not covered, and the 600 GeV resonance peak is captured by a single grid point (\textasciitilde{}3×10\textasciicircum{}4 fb vs \textasciitilde{}4×10\textasciicircum{}4 fb for the reference blue curve), understating the true peak height and width. The off-resonance minimum near 400 GeV sits at \textasciitilde{}9×10\textasciicircum{}2 fb versus \textasciitilde{}2×10\textasciicircum{}3 fb in the reference, a factor \textasciitilde{}2 low; the high-energy tail at 1200 GeV (\textasciitilde{}2.5×10\textasciicircum{}4 fb) agrees with the reference blue curve to within tens of percent.
\end{enumerate}
